% Bayesian Ranking -- Posterior rank averages (Laird & Louis 1989)
% Compile: pdflatex ranking.tex
\documentclass[varwidth=24cm,border=14pt]{standalone}
\usepackage[fleqn]{mathtools}
\usepackage{amssymb, bm}
\setlength{\mathindent}{0pt}
\begin{document}
\begin{minipage}{24cm}

\textbf{Posterior draws of regional baselines:}
\[
\alpha_j^s, \qquad j \in \{1,\dots,6\}\ \text{(regions)}, \quad s \in \{1,\dots,n\},\ n = 20{,}000
\]

\bigskip
\textbf{Step 1 -- Rank within each posterior sample:}

- Within each posterior draw, the six regional baselines $\alpha_j^s$ are ranked $1$--$6$.  

- Ties are assigned the mean of the ranks they share.

\bigskip
\textbf{Step 2 -- Average ranks across draws:}
\[
\hat{r}_j = \frac{1}{n} \sum_{s=1}^{n} r_j^s
\qquad \leftarrow \textbf{posterior mean rank}
\]

\bigskip
\textbf{Step 3 -- Rank the mean ranks:}
\[
r_j^{*} = \sum_{i=1}^{6} \mathbf{1}(\hat{r}_j > \hat{r}_i)
\qquad \leftarrow \textbf{final regional ranking}
\]

\end{minipage}
\end{document}
